\section{Conclusion and Future Work}
\label{sec:conclusion}

We addressed the architectural gap in biometric sensor integration for resource-constrained RTOSes. While ISO/IEC 19784 provides standardization for high-level systems and ISO/IEC 29164 specifies wire protocols, neither addresses RTOS-focused workflow, vendor-agnostic abstraction, type-safety, persistent state management, nor MCU-level memory constraints simultaneously. Recent work has explored various approaches to IoT biometric authentication\cite{abuhamad2021sensor,yang2021biometric}, but none addressed the specific constraints of multitasking RTOS environments. This forced developers into vendor-locked non-portable applications, application level state management, no type-safety, and blocking I/O.

We resolved these issues through BASE's three-plane architecture, which we validated through Zephyr RTOS. The design separates application API, subsystem contract, and driver implementation. This ensures type safety and normalized security policy. State management moved to driver-level, enforced by the subsystem. Section~\ref{sec:architecture} showed how this eliminates vendor-specific headers from application code. Section~\ref{sec:implementation} proved feasibility on the WeAct STM32F446RE with a minimum of \SI{3.5}{\kilo\byte} Flash and \SI{644}{\byte} RAM. Most importantly, Section~\ref{sec:evaluation} quantified portability by showing hardware migration required zero application code changes versus about 47 lines with legacy approaches. Security policy remains hardware-agnostic which makes it easy to swap the hardware without rewriting application code.

The hybrid interrupt-yielding I/O model addresses a fundamental RTOS requirement that blocking vendor libraries ignore. BASE yields the CPU during \num{200}--\SI{500}{\milli\second} sensor processing windows. By doing this it provides deterministic multitasking and power management. This is essential for industrial systems where biometric operations coexist with real-time control loops. We validated BASE across two physical sensor families directly on hardware (ZFM-X0, GT-5X) and confirmed modality-independent operation for camera-based multi-modal recognition (DFRobot AI10, face and palm vein) through datasheet compatibility analysis and an independent community implementation demonstrated in the Zephyr 4.4 public release\cite{zephyr_44_release}.

BASE's integration into Zephyr mainline as the official biometrics subsystem validates both technical feasibility and architectural necessity. This gives the IoT and industrial automation communities a standardized path to add biometric authentication without vendor lock-in while following RTOS conventions. As biometric authentication becomes mandatory for IoT security, BASE enables this transition on the resource-constrained edge devices that make up the majority of IoT deployments.

\subsection{Future Work}
\label{sec:future_work}

\textbf{FIDO2/WebAuthn support}: The FIDO2 standard enables passwordless authentication\cite{tran2024biometrics}. Extending BASE for edge FIDO2 authenticators requires a CTAP2 transport layer, binding of biometric user verification to credential operations, and secure key storage with attestation. This work is being pursued as a dedicated RTOS-level CTAP2 authenticator subsystem, with the BASE \texttt{biometric\_match()} result serving as the user verification primitive.